\documentclass[11pt]{article}
\usepackage[a4paper,margin=0.8in]{geometry}
\usepackage{calc}
\usepackage{amsmath,amssymb}
\usepackage{array,longtable,booktabs}
\usepackage{xcolor}
\usepackage{hyperref}
\usepackage{enumitem}
\usepackage{fancyvrb}
\usepackage{fvextra}
\usepackage{framed}
\usepackage{color}
\fvset{breaklines=true, breakanywhere=true}
\DefineVerbatimEnvironment{Highlighting}{Verbatim}{breaklines,breakanywhere,commandchars=\\\{\}}
\usepackage{color}
\usepackage{fancyvrb}
\newcommand{\VerbBar}{|}
\newcommand{\VERB}{\Verb[commandchars=\\\{\}]}
\DefineVerbatimEnvironment{Highlighting}{Verbatim}{commandchars=\\\{\}}
% Add ',fontsize=\small' for more characters per line
\usepackage{framed}
\definecolor{shadecolor}{RGB}{248,248,248}
\newenvironment{Shaded}{\begin{snugshade}}{\end{snugshade}}
\newcommand{\AlertTok}[1]{\textcolor[rgb]{0.94,0.16,0.16}{#1}}
\newcommand{\AnnotationTok}[1]{\textcolor[rgb]{0.56,0.35,0.01}{\textbf{\textit{#1}}}}
\newcommand{\AttributeTok}[1]{\textcolor[rgb]{0.13,0.29,0.53}{#1}}
\newcommand{\BaseNTok}[1]{\textcolor[rgb]{0.00,0.00,0.81}{#1}}
\newcommand{\BuiltInTok}[1]{#1}
\newcommand{\CharTok}[1]{\textcolor[rgb]{0.31,0.60,0.02}{#1}}
\newcommand{\CommentTok}[1]{\textcolor[rgb]{0.56,0.35,0.01}{\textit{#1}}}
\newcommand{\CommentVarTok}[1]{\textcolor[rgb]{0.56,0.35,0.01}{\textbf{\textit{#1}}}}
\newcommand{\ConstantTok}[1]{\textcolor[rgb]{0.56,0.35,0.01}{#1}}
\newcommand{\ControlFlowTok}[1]{\textcolor[rgb]{0.13,0.29,0.53}{\textbf{#1}}}
\newcommand{\DataTypeTok}[1]{\textcolor[rgb]{0.13,0.29,0.53}{#1}}
\newcommand{\DecValTok}[1]{\textcolor[rgb]{0.00,0.00,0.81}{#1}}
\newcommand{\DocumentationTok}[1]{\textcolor[rgb]{0.56,0.35,0.01}{\textbf{\textit{#1}}}}
\newcommand{\ErrorTok}[1]{\textcolor[rgb]{0.64,0.00,0.00}{\textbf{#1}}}
\newcommand{\ExtensionTok}[1]{#1}
\newcommand{\FloatTok}[1]{\textcolor[rgb]{0.00,0.00,0.81}{#1}}
\newcommand{\FunctionTok}[1]{\textcolor[rgb]{0.13,0.29,0.53}{\textbf{#1}}}
\newcommand{\ImportTok}[1]{#1}
\newcommand{\InformationTok}[1]{\textcolor[rgb]{0.56,0.35,0.01}{\textbf{\textit{#1}}}}
\newcommand{\KeywordTok}[1]{\textcolor[rgb]{0.13,0.29,0.53}{\textbf{#1}}}
\newcommand{\NormalTok}[1]{#1}
\newcommand{\OperatorTok}[1]{\textcolor[rgb]{0.81,0.36,0.00}{\textbf{#1}}}
\newcommand{\OtherTok}[1]{\textcolor[rgb]{0.56,0.35,0.01}{#1}}
\newcommand{\PreprocessorTok}[1]{\textcolor[rgb]{0.56,0.35,0.01}{\textit{#1}}}
\newcommand{\RegionMarkerTok}[1]{#1}
\newcommand{\SpecialCharTok}[1]{\textcolor[rgb]{0.81,0.36,0.00}{\textbf{#1}}}
\newcommand{\SpecialStringTok}[1]{\textcolor[rgb]{0.31,0.60,0.02}{#1}}
\newcommand{\StringTok}[1]{\textcolor[rgb]{0.31,0.60,0.02}{#1}}
\newcommand{\VariableTok}[1]{\textcolor[rgb]{0.00,0.00,0.00}{#1}}
\newcommand{\VerbatimStringTok}[1]{\textcolor[rgb]{0.31,0.60,0.02}{#1}}
\newcommand{\WarningTok}[1]{\textcolor[rgb]{0.56,0.35,0.01}{\textbf{\textit{#1}}}}
\setlist[itemize]{leftmargin=*}
\setlist[enumerate]{leftmargin=*}
\hypersetup{colorlinks=true,linkcolor=blue,urlcolor=blue}
\providecommand{\tightlist}{%
  \setlength{\itemsep}{0pt}\setlength{\parskip}{0pt}}
\title{R Workshop -- Day 7 -- Reproducible Research and Capstone Presentation}
\author{Applied R for Statistical Teaching, Research and Reproducible Analysis}
\date{}
\begin{document}
\maketitle
\tableofcontents
\newpage
\hypertarget{r-workshop-day-7-reproducible-research-and-capstone-presentation}{%
\section{R Workshop -- Day 7 -- Reproducible Research and Capstone
Presentation}\label{r-workshop-day-7-reproducible-research-and-capstone-presentation}}

\textbf{Session 7 of 7 \textbar{} Duration: 3 hours} \textbf{Programme:}
Applied R for Statistical Teaching, Research and Reproducible Analysis

\hypertarget{learning-objectives}{%
\subsection{Learning objectives}\label{learning-objectives}}

By the end of this session you will be able to:

\begin{enumerate}
\def\labelenumi{\arabic{enumi}.}
\tightlist
\item
  Organise a project into raw data, processed data, scripts, outputs and
  reports, documented with a README.
\item
  Build a Quarto report with YAML metadata, code-cell options and inline
  results.
\item
  Render the same report to HTML, Word and PDF.
\item
  Test a report in a clean R session and record session information.
\item
  Conduct a structured peer review using a short checklist.
\item
  Present a capstone analysis clearly in five minutes.
\end{enumerate}

\hypertarget{capstone-link}{%
\subsection{Capstone link}\label{capstone-link}}

Today integrates everything from Days 2-6 -- the cleaned dataset, the
graphics, the inferential result and the diagnosed model -- into one
reproducible report, and closes the workshop with the Day 1 baseline
task repeated for comparison.

\hypertarget{segment-plan-3-hours}{%
\subsection{Segment plan (3 hours)}\label{segment-plan-3-hours}}

\begin{longtable}[]{@{}
  >{\raggedright\arraybackslash}p{(\columnwidth - 4\tabcolsep) * \real{0.3000}}
  >{\raggedleft\arraybackslash}p{(\columnwidth - 4\tabcolsep) * \real{0.4000}}
  >{\raggedright\arraybackslash}p{(\columnwidth - 4\tabcolsep) * \real{0.3000}}@{}}
\toprule\noalign{}
\begin{minipage}[b]{\linewidth}\raggedright
Segment
\end{minipage} & \begin{minipage}[b]{\linewidth}\raggedleft
Time
\end{minipage} & \begin{minipage}[b]{\linewidth}\raggedright
Purpose
\end{minipage} \\
\midrule\noalign{}
\endhead
\bottomrule\noalign{}
\endlastfoot
Opening and recap & 10 min & Recap the full capstone thread, Day 2
through Day 6 \\
Concept bridge & 25 min & Why reproducibility is a deliverable, not an
afterthought \\
Guided coding & 60 min & Build and render the Quarto capstone report \\
Participant practice & 45 min & Integrate your own Days 2-6 artefacts
into the template \\
Interpretation and reporting & 25 min & Peer review using the
checklist \\
Evidence log and capstone update & 15 min & Repeat the Day 1 baseline
task and present \\
\end{longtable}

\begin{center}\rule{0.5\linewidth}{0.5pt}\end{center}

\hypertarget{part-1-concept-bridge-core-workshop-activity}{%
\subsection{Part 1 -- Concept bridge (Core workshop
activity)}\label{part-1-concept-bridge-core-workshop-activity}}

\hypertarget{why-reproducibility-is-a-deliverable}{%
\subsubsection{Why reproducibility is a
deliverable}\label{why-reproducibility-is-a-deliverable}}

A result that only runs on your machine, in your current R session, with
objects you created interactively three days ago, is not yet a finished
piece of work. Reproducibility means: a colleague (or you, in six
months) can open the project, run one command, and get the same report
-- numbers, figures and all -- without manual intervention.

\hypertarget{project-organisation}{%
\subsubsection{Project organisation}\label{project-organisation}}

\begin{Shaded}
\begin{Highlighting}[]
\NormalTok{your{-}capstone{-}project/}
\NormalTok{|{-}{-} your{-}capstone{-}project.Rproj}
\NormalTok{|{-}{-} Data/}
\NormalTok{|   |{-}{-} raw/}
\NormalTok{|   \textasciigrave{}{-}{-} processed/}
\NormalTok{|{-}{-} Scripts/}
\NormalTok{|{-}{-} Outputs/}
\NormalTok{|{-}{-} Reports/}
\NormalTok{|   \textasciigrave{}{-}{-} capstone{-}report.qmd}
\NormalTok{\textasciigrave{}{-}{-} README.md}
\end{Highlighting}
\end{Shaded}

Relative paths (\texttt{Data/raw/anchor\_dataset.csv}, not
\texttt{C:/Users/you/Desktop/...}) are what make this folder portable
between machines -- the same principle from Day 1, now applied to a
finished report.

\hypertarget{quarto-basics}{%
\subsubsection{Quarto basics}\label{quarto-basics}}

A \texttt{.qmd} file has a YAML header (metadata: title, author, output
formats) followed by Markdown text interleaved with executable code
cells. Quarto is the modern successor to R Markdown and renders the same
source to multiple output formats without changing the document.

\begin{center}\rule{0.5\linewidth}{0.5pt}\end{center}

\hypertarget{part-2-guided-coding-core-workshop-activity}{%
\subsection{Part 2 -- Guided coding (Core workshop
activity)}\label{part-2-guided-coding-core-workshop-activity}}

Open \texttt{Quarto-Template/capstone-report.qmd}.

\hypertarget{yaml-metadata}{%
\subsubsection{2.1 YAML metadata}\label{yaml-metadata}}

\begin{Shaded}
\begin{Highlighting}[]
\PreprocessorTok{{-}{-}{-}}
\FunctionTok{title}\KeywordTok{:}\AttributeTok{ }\StringTok{"R Workshop Capstone Report"}
\FunctionTok{author}\KeywordTok{:}\AttributeTok{ }\StringTok{"Participant name"}
\FunctionTok{format}\KeywordTok{:}
\AttributeTok{  }\FunctionTok{html}\KeywordTok{:}\AttributeTok{ default}
\AttributeTok{  }\FunctionTok{docx}\KeywordTok{:}\AttributeTok{ default}
\AttributeTok{  }\FunctionTok{pdf}\KeywordTok{:}\AttributeTok{ default}
\FunctionTok{execute}\KeywordTok{:}
\AttributeTok{  }\FunctionTok{echo}\KeywordTok{:}\AttributeTok{ }\CharTok{true}
\AttributeTok{  }\FunctionTok{warning}\KeywordTok{:}\AttributeTok{ }\CharTok{false}
\AttributeTok{  }\FunctionTok{message}\KeywordTok{:}\AttributeTok{ }\CharTok{false}
\PreprocessorTok{{-}{-}{-}}
\end{Highlighting}
\end{Shaded}

\texttt{format:} lists every output you want to render to.
\texttt{execute:} options apply to every code cell unless overridden
per-cell.

\hypertarget{code-cell-options}{%
\subsubsection{2.2 Code-cell options}\label{code-cell-options}}

\begin{Shaded}
\begin{Highlighting}[]
\CommentTok{\#| label: setup}
\CommentTok{\#| echo: false}
\FunctionTok{library}\NormalTok{(tidyverse)}
\end{Highlighting}
\end{Shaded}

\texttt{\#\textbar{}\ label:} names the cell (useful for
cross-references and error messages);
\texttt{\#\textbar{}\ echo:\ false} hides the code but still runs it and
shows any output; \texttt{\#\textbar{}\ fig-cap:} adds a caption to a
plot produced in that cell.

\hypertarget{figures-tables-equations-and-inline-results}{%
\subsubsection{2.3 Figures, tables, equations and inline
results}\label{figures-tables-equations-and-inline-results}}

\begin{Shaded}
\begin{Highlighting}[]
\CommentTok{\#| label: fig{-}score{-}by{-}track}
\CommentTok{\#| fig{-}cap: "Post{-}training score by training track"}
\NormalTok{lecturers }\OtherTok{\textless{}{-}} \FunctionTok{read\_csv}\NormalTok{(}\StringTok{"Data/processed/anchor\_dataset\_clean.csv"}\NormalTok{, }\AttributeTok{show\_col\_types =} \ConstantTok{FALSE}\NormalTok{)}
\FunctionTok{ggplot}\NormalTok{(lecturers, }\FunctionTok{aes}\NormalTok{(}\AttributeTok{x =}\NormalTok{ training\_track, }\AttributeTok{y =}\NormalTok{ post\_score)) }\SpecialCharTok{+}
  \FunctionTok{geom\_boxplot}\NormalTok{(}\AttributeTok{fill =} \StringTok{"\#FEB24C"}\NormalTok{) }\SpecialCharTok{+}
  \FunctionTok{theme\_minimal}\NormalTok{()}
\end{Highlighting}
\end{Shaded}

Inline results, e.g.~``the mean attendance was
\texttt{\textasciigrave{}r\ round(mean(lecturers\$attendance,\ na.rm\ =\ TRUE),\ 1)\textasciigrave{}}
percent'', insert a live-computed number directly into the prose -- so
the number updates automatically if the data changes.

\hypertarget{rendering}{%
\subsubsection{2.4 Rendering}\label{rendering}}

\begin{Shaded}
\begin{Highlighting}[]
\NormalTok{quarto}\SpecialCharTok{::}\FunctionTok{quarto\_render}\NormalTok{(}\StringTok{"Reports/capstone{-}report.qmd"}\NormalTok{)}
\end{Highlighting}
\end{Shaded}

or, from the terminal:
\texttt{quarto\ render\ "Reports/capstone-report.qmd"}. This produces
\texttt{.html}, \texttt{.docx} and \texttt{.pdf} versions in one step,
all from the same source.

\hypertarget{testing-in-a-clean-session-and-recording-sessioninfo}{%
\subsubsection{\texorpdfstring{2.5 Testing in a clean session and
recording
\texttt{sessionInfo()}}{2.5 Testing in a clean session and recording sessionInfo()}}\label{testing-in-a-clean-session-and-recording-sessioninfo}}

Before considering a report finished, restart R
(\texttt{Session\ -\textgreater{}\ Restart\ R} in RStudio) and re-render
from scratch -- this catches any hidden dependency on an object that
only exists because you ran something manually earlier. Always include:

\begin{Shaded}
\begin{Highlighting}[]
\CommentTok{\#| label: session{-}info}
\FunctionTok{sessionInfo}\NormalTok{()}
\end{Highlighting}
\end{Shaded}

near the end of the report, so the exact package versions used are
recorded alongside the results.

\hypertarget{peer-review-checklist}{%
\subsubsection{2.6 Peer review checklist}\label{peer-review-checklist}}

Exchange reports with a partner and check:

\begin{itemize}
\tightlist
\item[$\square$]
  Does the report render from a clean session without errors?
\item[$\square$]
  Is the research question stated clearly in the first section?
\item[$\square$]
  Are raw and processed data kept separate, with cleaning decisions
  documented?
\item[$\square$]
  Does every figure have a title, labelled axes with units, and a
  caption?
\item[$\square$]
  Is the inferential or model result reported with both a number and a
  plain-language sentence?
\item[$\square$]
  Is \texttt{sessionInfo()} included?
\end{itemize}

\begin{center}\rule{0.5\linewidth}{0.5pt}\end{center}

\hypertarget{part-3-participant-practice-core-workshop-activity}{%
\subsection{Part 3 -- Participant practice (Core workshop
activity)}\label{part-3-participant-practice-core-workshop-activity}}

\begin{enumerate}
\def\labelenumi{\arabic{enumi}.}
\tightlist
\item
  Copy your Day 2 cleaning script's output, your three Day 3 graphics,
  your Day 4 inferential result and your Day 5/6 model into
  \texttt{Quarto-Template/capstone-report.qmd}, replacing the
  placeholder text in each section.
\item
  Render to HTML, Word and PDF (or note in the report if a particular
  renderer is unavailable on your machine).
\item
  Restart R and re-render once more to confirm the report is
  reproducible from a clean session.
\item
  Exchange reports with a partner and complete the peer-review checklist
  for each other's report.
\item
  Repeat the Day 1 baseline task on
  \texttt{Data/raw/day7\_baseline\_sample.csv} and add a short final
  section comparing your Day 1 and Day 7 work.
\item
  Prepare a five-minute presentation: research question, one figure, one
  result, one limitation.
\end{enumerate}

A facilitator-reviewed, fully completed example report is in
\texttt{Solution-Scripts/day7\_capstone\_report\_solution.qmd}.

\begin{center}\rule{0.5\linewidth}{0.5pt}\end{center}

\hypertarget{optional-demonstration}{%
\subsection{Optional demonstration}\label{optional-demonstration}}

\begin{itemize}
\tightlist
\item
  \textbf{Confidential-data handling} -- never commit raw
  participant-identifiable data to a shared repository; demonstrate
  \texttt{.gitignore} patterns and anonymisation as a brief discussion,
  not a hands-on exercise.
\item
  \textbf{Multi-format formatting differences} -- show how the same
  Quarto source renders slightly differently in HTML versus PDF
  (e.g.~figure sizing, table wrapping) and how to use format-specific
  code-cell options to handle this.
\item
  \textbf{Advanced Quarto options} -- cross-references
  (\texttt{@fig-...}), citations (\texttt{.bib} files), parameterised
  reports.
\end{itemize}

\begin{center}\rule{0.5\linewidth}{0.5pt}\end{center}

\hypertarget{reference-or-take-home-material}{%
\subsection{Reference or take-home
material}\label{reference-or-take-home-material}}

\begin{itemize}
\tightlist
\item
  Troubleshooting rendering problems: missing LaTeX packages for PDF
  output, missing pandoc for Word output, encoding issues with special
  characters.
\item
  A report-polishing checklist for after the workshop: consistent figure
  styling, a finalised abstract/summary, spell-checked prose.
\end{itemize}

\begin{center}\rule{0.5\linewidth}{0.5pt}\end{center}

\hypertarget{daily-deliverable}{%
\subsection{Daily deliverable}\label{daily-deliverable}}

A rendered, reproducible capstone report (HTML, Word and/or PDF) and a
five-minute presentation.

\hypertarget{evidence-log-entry}{%
\subsubsection{Evidence log entry}\label{evidence-log-entry}}

Complete the Day 7 row in \texttt{Workshop-Evidence-Log.md}. Fill in the
pre/post assessment task comparison: how did your Day 7 baseline-task
script differ from your Day 1 version, and what does that difference
show about your progress through the workshop?

\begin{center}\rule{0.5\linewidth}{0.5pt}\end{center}

\hypertarget{facilitator-notes}{%
\subsection{Facilitator notes}\label{facilitator-notes}}

\begin{itemize}
\tightlist
\item
  Day 7 is for integration and presentation, not for starting new
  analysis -- if a participant has not completed earlier days'
  deliverables, help them integrate what exists rather than trying to
  backfill missing work live.
\item
  PDF rendering depends on a working LaTeX installation; confirm this in
  advance, and have HTML as a guaranteed fallback format for every
  participant.
\item
  Keep presentations strictly timed at five minutes; this is itself a
  transferable skill (communicating a result to a non-specialist
  audience under a hard time constraint), not just a logistics choice.
\end{itemize}
\end{document}
